\appendix

\section{Supplementary diagnostics and reproducibility evidence}\label{sec:appendix}

This appendix contains the current evidence-bearing diagnostics used to
interpret the weak-label baseline. It is limited to reproducible feature,
surface-model, canopy, and spatial-fold evidence. Retired
external-triangulation diagnostics, duplicate weak-label panels, manual-audit
figures, and untrained deep-model outputs are not part of the supplementary
bundle; the completed audit is reported in the main text.

All panels were generated from the locked feature matrix and model-output
directories identified in the project figure map. Source paths, parameters,
selection rules, and claim boundaries are retained in the accompanying
manifests and build scripts.

\subsection{Matched urban-texture, surface-model, and canopy atlas}

The atlas contains the same 48 deterministic selection cells used for the
within-class climate--morphology comparison. Each cell is shown through
matched figure-ground texture, surface-model elevation/roughness, and cleaned
10 m canopy-height views. The atlas is descriptive and reproducible from the
exported selection manifest; it is not an audited local subtype validation.

\begin{figure*}[htbp]
    \centering
    \includegraphics[width=\textwidth]{fig_s1_texture_dsm_canopy_atlas_48_cells.png}
    \caption{Matched 48-cell urban-texture atlas. Each selected cell contains
    figure-ground texture, surface-model elevation/roughness, and canopy-height
    views on shared robust scales. The figure is a visual evidence bridge, not
    a validated building-height or audited local-classification result.}
    \label{fig:s1-texture-surface-canopy}
\end{figure*}
\FloatBarrier

\subsection{Feature quality and collinearity}

Feature missingness is interpreted together with structural-absence flags
before fold-specific imputation. The collinearity panel records the largest
groups at absolute correlation $r\geq0.95$; family-level ablations and
representative variables are used to reduce redundancy in the reported
baselines.

\begin{figure*}[htbp]
    \centering
    \includegraphics[width=0.96\textwidth]{fig_s2_feature_quality.png}
    \caption{Feature-quality diagnostic. (a) Mean and maximum missingness by
    family, (b) feature count versus near-constant columns, and (c) the largest
    high-correlation feature groups.}
    \label{fig:s2-feature-quality}
\end{figure*}
\FloatBarrier

\subsection{Spatial-fold diagnostics}

The primary spatial folds are fixed before model comparison. This diagnostic
shows class support by fold and the fold-wise distribution of the primary
LightGBM out-of-fold prediction probability. It is a stability diagnostic for
the weak-label baseline and does not constitute audited subtype validation.

\begin{figure*}[htbp]
    \centering
    \includegraphics[width=0.92\textwidth]{fig_s3_spatial_fold_diagnostics.png}
    \caption{Spatial-fold diagnostics. (a) Weak-class support across the fixed
    spatial folds and (b) primary LightGBM out-of-fold prediction probability
    by fold.}
    \label{fig:s3-spatial-folds}
\end{figure*}
\FloatBarrier

\subsection{Reproducibility checklist}

The analysis package should be considered internally reproducible when the
following checks pass: (i) the v8 input has 16,506 rows and 374 fields; (ii)
all manifest-referenced fields exist; (iii) the four-class filter returns
5,339 rows; (iv) every eligible row has one spatial fold; (v) preprocessing is
fitted inside training folds; (vi) all reported out-of-fold outputs are
regenerated from the locked scripts and parameters; and (vii) the PDF compiles
without missing figure, table, citation, or reference warnings. The manual
audit is complete, but deep multimodal fusion remains unavailable until its
output manifests exist, so no supplementary panel should be read as evidence
of a trained deep model.
